\documentclass{situationist-preprint}

\title{The Question of Uptake}
\msid{AES-15-b}
\msfamily{AES-15}
\msdiscipline{Aesthetics}
\mstype{Companion manuscript}
\mskeywords{machine art; institutional theory; critical uptake; art market}

\begin{document}
\maketitle
\vspace{-1em}

\begin{abstract}
AES-15-a argues that machine outputs can be art if an artworld takes them up, and that whether it does is open (Proposition~3.2). We make the openness precise. Critical uptake and market uptake are distinct, the second is not evidence of the first, and what would close the question is a pattern of critical practice, not a sale or a philosophical argument.
\end{abstract}

\section{Two kinds of uptake}

By Definition~2.2 of AES-15-a, an object receives critical uptake when the artworld's critical practice engages it as a candidate: it is interpreted, compared with other works and assessed by the standards applied to artworks. Market uptake is the treatment of an object as a good with a price. The two often coincide for artworks, but they come apart in both directions: works neglected by the market are discussed by critics, and collectibles, souvenirs and memorabilia are priced without being interpreted.

\begin{proposition}\label{p:distinct}
Market uptake of a machine output is neither necessary nor sufficient for its critical uptake, and evidence of the first is weak evidence of the second.
\end{proposition}

The weakness of the evidence has a specific source. A price can be paid for an object as a work, as a document of a technology, or as a novelty, and the transaction does not reveal which. Critical uptake, by contrast, is revealed by its products: interpretations that attribute meaning to a particular output, comparisons that place it among other works, and judgments of success or failure that appeal to standards of art rather than of engineering \citep{Danto1964,Dickie1974}.

\section{What would close the question}

The question of Proposition~3.2 of AES-15-a would be closed in the affirmative by a stable practice in which particular machine outputs are interpreted and assessed as works, and in the negative by a practice that persistently discusses them only as products of a technique. At present criticism does both, often of the same object. Philosophy can state the conditions; their satisfaction is a matter of what critics come to do, which is why the question remains open.

\begin{thebibliography}{9}\small
\bibitem[Danto(1964)]{Danto1964} Danto, A.~C. (1964). The artworld. \emph{The Journal of Philosophy}, 61(19), 571--584.
\bibitem[Dickie(1974)]{Dickie1974} Dickie, G. (1974). \emph{Art and the Aesthetic: An Institutional Analysis}. Cornell University Press.
\end{thebibliography}

\end{document}
